\documentclass{article}
\usepackage[utf8]{inputenc}
\usepackage{hyperref}
\usepackage{booktabs}

\title{PURE and CORE Draft Index}
\author{Research Notes}
\date{2026}

\begin{document}
\maketitle

\section*{Draft Index}
This directory contains working notes for the current architecture and report direction. Report-facing terminology should use \textbf{PURE} for \emph{Predicate-aware Uncropped Relation Embedding} and \textbf{CORE} for \emph{Compositional Relation Evaluation}.

\begin{itemize}
  \item \textbf{PURE note:} \texttt{notes/paper1\_pure.tex} --- current architecture, phase curriculum, three training pillars, and reporting rules.
  \item \textbf{CORE note:} \texttt{notes/paper2\_core.tex} --- dataset status, six-group diagnostic role, and future evaluation plan.
  \item \textbf{Diagrams:} \texttt{notes/diagrams.tex} --- figures aligned with the current PURE code path.
  \item \textbf{Current status:} \texttt{notes/current\_status.tex} --- concise implementation snapshot and report guidance.
\end{itemize}

\section*{Do Not Overclaim}
Branch-ramp runs are debug/ablation runs. CORE exists as generated data, but it is not yet the main training or evaluation protocol for PURE. Final report claims should separate no-prior metrics, calibration metrics, and future-work items.

\section*{Build Notes}
Compile drafts independently from the repository root with \texttt{pdflatex -interaction=nonstopmode -output-directory=notes notes/<file>.tex}.

\end{document}